\documentclass[12pt,a4paper]{article}
\usepackage[utf8]{inputenc}
\usepackage[T1]{fontenc}
\usepackage[spanish]{babel}
\usepackage{geometry}
\usepackage{graphicx}
\usepackage{hyperref}
\usepackage{xcolor}
\usepackage{titlesec}
\usepackage{enumitem}
\usepackage{fancyhdr}
\usepackage{booktabs}
\usepackage{tabularx}
\usepackage{listings}
\usepackage{tikz}
\usetikzlibrary{shapes.geometric, arrows.meta, positioning}

\geometry{margin=2.5cm}

% Colores institucionales
\definecolor{primary}{RGB}{106, 27, 154}
\definecolor{secondary}{RGB}{243, 232, 255}
\definecolor{codebg}{RGB}{245, 245, 245}
\definecolor{linkcolor}{RGB}{106, 27, 154}

% Hipervínculos
\hypersetup{
    colorlinks=true,
    linkcolor=primary,
    urlcolor=linkcolor,
    citecolor=primary
}

% Estilo de código
\lstset{
    backgroundcolor=\color{codebg},
    basicstyle=\ttfamily\footnotesize,
    breaklines=true,
    frame=single,
    rulecolor=\color{primary!30},
    keywordstyle=\color{primary}\bfseries,
    commentstyle=\color{gray}\itshape,
    numbers=left,
    numberstyle=\tiny\color{gray},
    numbersep=5pt,
    tabsize=4
}

% Encabezado y pie de página
\pagestyle{fancy}
\fancyhf{}
\fancyhead[L]{\textcolor{primary}{\textbf{Vialert}}}
\fancyhead[R]{\textcolor{gray}{\small Proyecto Final --- Fase 3}}
\fancyfoot[C]{\thepage}
\renewcommand{\headrulewidth}{0.4pt}
\renewcommand{\headrule}{\hbox to\headwidth{\color{primary}\leaders\hrule height \headrulewidth\hfill}}

% Títulos de sección
\titleformat{\section}{\large\bfseries\color{primary}}{}{0em}{}[\titlerule]
\titleformat{\subsection}{\normalsize\bfseries\color{primary!80}}{}{0em}{}

% ─────────────────────────────────────────────
\begin{document}

% ── PORTADA ──────────────────────────────────
\begin{titlepage}
    \centering
    \vspace*{2cm}
    {\huge\bfseries\color{primary} Vialert\par}
    \vspace{0.4cm}
    {\large\color{gray} Aplicación Ciudadana de Reportes Viales\par}
    \vspace{1.5cm}
    \rule{\textwidth}{1.5pt}\par
    \vspace{0.3cm}
    {\LARGE\bfseries Entrega Final --- Fase 3\par}
    \vspace{0.3cm}
    \rule{\textwidth}{1.5pt}\par
    \vspace{1.5cm}
    {\large\textbf{Temática 1: Reportes ciudadanos}\par}
    \vspace{0.5cm}
    {\large Universidad del Quindío\par}
    {\large Ingeniería de Software / Desarrollo Móvil\par}
    \vspace{1.5cm}
    {\large\textbf{Autor:} Jordy Díaz\par}
    \vspace{0.3cm}
    {\large\textbf{Correo:} JFDIAZB@uqvirtual.edu.co\par}
    \vfill
    {\large Mayo 2026\par}
\end{titlepage}

% ── TABLA DE CONTENIDOS ───────────────────────
\tableofcontents
\newpage

% ════════════════════════════════════════════
\section{Repositorio del Proyecto}
% ════════════════════════════════════════════

El código fuente completo de la aplicación Vialert se encuentra disponible en el siguiente repositorio de GitHub:

\begin{center}
    \colorbox{secondary}{\parbox{0.85\textwidth}{\centering\vspace{6pt}
        \url{https://github.com/JordyDiazB/Vialert}
    \vspace{6pt}}}
\end{center}

El repositorio contiene el proyecto Android completo desarrollado con Kotlin y Jetpack Compose, incluyendo todas las integraciones de servicios en la nube descritas en este documento.

% ════════════════════════════════════════════
\section{Descripción General de la Aplicación}
% ════════════════════════════════════════════

\textbf{Vialert} es una aplicación móvil Android orientada a la participación ciudadana en Colombia, específicamente diseñada para reportar y gestionar problemas de infraestructura vial, alumbrado público y seguridad vial. La plataforma conecta a ciudadanos con moderadores administrativos mediante un flujo completo de reporte, moderación y resolución.

\subsection{Características principales}
\begin{itemize}[leftmargin=1.5em]
    \item Registro e inicio de sesión con Firebase Authentication
    \item Creación de reportes geolocalizados con foto, categoría y descripción
    \item Feed de reportes con filtro por distancia usando GPS real
    \item Mapa interactivo de reportes (Mapbox)
    \item Sistema de moderación con análisis automático por Inteligencia Artificial
    \item Asistente virtual conversacional integrado
    \item Sistema de reputación ciudadana con 4 niveles
    \item Notificaciones push al verificar o rechazar un reporte
    \item Panel de administración con estadísticas y moderación
\end{itemize}

% ════════════════════════════════════════════
\section{Persistencia de Datos: Firebase Firestore con Caché}
% ════════════════════════════════════════════

\subsection{Implementación de Firestore}

La aplicación utiliza \textbf{Firebase Cloud Firestore} como base de datos principal en tiempo real. Los datos se organizan en las colecciones \texttt{users} y \texttt{reports}, con listeners en tiempo real que actualizan la UI automáticamente mediante \texttt{StateFlow}.

\subsection{Caché offline persistente}

Se habilitó la persistencia offline de Firestore mediante \texttt{PersistentCacheSettings}, lo que permite que la aplicación funcione con datos previamente descargados cuando no hay conexión a internet:

\begin{lstlisting}[language=Kotlin, caption={Configuración de caché persistente en FirebaseModule.kt}]
val cacheSettings = PersistentCacheSettings.newBuilder()
    .setSizeBytes(FirebaseFirestoreSettings.CACHE_SIZE_UNLIMITED)
    .build()
val settings = FirebaseFirestoreSettings.Builder()
    .setLocalCacheSettings(cacheSettings)
    .build()
FirebaseFirestore.getInstance().apply { firestoreSettings = settings }
\end{lstlisting}

Con esta configuración, Firestore almacena localmente en el dispositivo todos los documentos consultados, permitiendo lecturas sin conexión y sincronización automática al recuperar la conectividad.

% ════════════════════════════════════════════
\section{Autenticación con Firebase}
% ════════════════════════════════════════════

La autenticación de usuarios se gestiona completamente con \textbf{Firebase Authentication}, implementando:

\begin{itemize}[leftmargin=1.5em]
    \item \textbf{Registro} de nuevos usuarios con correo electrónico y contraseña
    \item \textbf{Inicio de sesión} con validación de credenciales
    \item \textbf{Recuperación de contraseña} mediante envío de correo de restablecimiento
    \item \textbf{Eliminación de cuenta} (baja del usuario desde Firebase Auth y Firestore simultáneamente)
    \item \textbf{Roles de usuario}: ciudadano y administrador, gestionados en Firestore
\end{itemize}

\subsection{Recuperación de contraseña}

La funcionalidad de recuperación envía un correo oficial de Firebase al usuario con un enlace para restablecer su contraseña:

\begin{lstlisting}[language=Kotlin, caption={Envío de correo de recuperación}]
override suspend fun sendPasswordReset(email: String) {
    auth.sendPasswordResetEmail(email).await()
}
\end{lstlisting}

El flujo completo incluye una pantalla dedicada (\texttt{ForgetPasswordScreen}) con validación del correo y mensajes de confirmación al usuario.

% ════════════════════════════════════════════
\section{Funcionalidad de Mapas: Mapbox}
% ════════════════════════════════════════════

La aplicación integra el \textbf{SDK de Mapbox Maps v11} para Android en dos contextos:

\begin{enumerate}[leftmargin=1.5em]
    \item \textbf{Exploración de reportes}: pantalla de mapa interactivo con marcadores de todos los reportes activos, agrupados por categoría y con indicadores de severidad.
    \item \textbf{Creación de reportes}: selector de ubicación que permite al usuario tocar el mapa para establecer las coordenadas exactas del problema reportado.
\end{enumerate}

\subsection{Filtro por distancia con GPS}

El feed de reportes incluye filtros de proximidad que usan la ubicación real del dispositivo mediante \texttt{FusedLocationProviderClient} y la fórmula de Haversine para calcular la distancia entre coordenadas:

\begin{itemize}[leftmargin=1.5em]
    \item \textbf{Todos}: sin filtro de distancia
    \item \textbf{Cercanos}: reportes en radio de 5 km
    \item \textbf{En la ciudad}: reportes en radio de 30 km
\end{itemize}

% ════════════════════════════════════════════
\section{Almacenamiento de Imágenes en la Nube: Cloudinary}
% ════════════════════════════════════════════

\subsection{Servicio utilizado}

Se utilizó \textbf{Cloudinary} como servicio de almacenamiento y entrega de imágenes en la nube. Cloudinary es una plataforma especializada en gestión de activos multimedia con CDN global, transformaciones automáticas y URLs de acceso seguro (HTTPS).

\subsection{Implementación}

Las imágenes se suben directamente desde el dispositivo al endpoint de Cloudinary mediante una petición HTTP multipart, utilizando un \textit{upload preset} sin firma (\textit{unsigned upload}):

\begin{lstlisting}[language=Kotlin, caption={Subida de imágenes a Cloudinary (ImageUploadService.kt)}]
private val cloudName = "dupm8asvi"
private val uploadPreset = "vialert_preset"

suspend fun uploadImage(uri: Uri, context: Context, folder: String): String {
    val url = "https://api.cloudinary.com/v1_1/$cloudName/image/upload"
    // Multipart/form-data con file, upload_preset y folder
    // Retorna la secure_url del CDN
}
\end{lstlisting}

\subsection{Uso en la aplicación}

\begin{itemize}[leftmargin=1.5em]
    \item \textbf{Reportes}: las fotos adjuntadas al crear un reporte se suben a la carpeta \texttt{reports/} y la URL se guarda en Firestore.
    \item \textbf{Perfil de usuario}: la foto de perfil se sube a la carpeta \texttt{profiles/} al editar el perfil.
\end{itemize}

Las imágenes se muestran en la app usando la librería \textbf{Coil} con carga asíncrona y caché local.

% ════════════════════════════════════════════
\section{Implementación de Inteligencia Artificial}
% ════════════════════════════════════════════

\subsection{Descripción funcional}

La inteligencia artificial se integra en dos funcionalidades principales:

\begin{enumerate}[leftmargin=1.5em]
    \item \textbf{Moderación automática de reportes}: el sistema analiza cada reporte ciudadano y genera una recomendación (VERIFICAR / RECHAZAR) con nivel de confianza, puntos clave y justificación. Esto asiste al administrador en la toma de decisiones.
    \item \textbf{Asistente virtual conversacional}: un chatbot integrado en la pantalla principal que responde preguntas de los ciudadanos sobre el uso de la app, categorías de reportes, seguridad vial e infraestructura.
\end{enumerate}

\subsection{Servicio de IA utilizado}

Se integró la API de \textbf{OpenRouter} (\url{https://openrouter.ai}), que ofrece acceso unificado a múltiples modelos de lenguaje grande (LLMs) a través de una interfaz compatible con OpenAI. OpenRouter permite utilizar modelos gratuitos de proveedores como Meta, Google y NVIDIA sin costo por token.

\textbf{Modelos utilizados} (en orden de preferencia con fallback automático):
\begin{itemize}[leftmargin=1.5em]
    \item \texttt{meta-llama/llama-3.3-70b-instruct:free} (principal)
    \item \texttt{google/gemma-4-26b-a4b-it:free} (secundario)
    \item \texttt{meta-llama/llama-3.2-3b-instruct:free} (fallback)
\end{itemize}

\subsection{Diagrama de integración}

\begin{center}
\begin{tikzpicture}[
    node distance=1.4cm and 2.2cm,
    box/.style={rectangle, rounded corners=6pt, draw=primary, fill=secondary,
                minimum width=2.8cm, minimum height=1cm, text centered, font=\small},
    cloud/.style={ellipse, draw=primary!60, fill=primary!10,
                  minimum width=3cm, minimum height=1cm, text centered, font=\small},
    arr/.style={-{Stealth[length=6pt]}, thick, color=primary!70}
]

\node[box] (vm)       {Admin / User\\ViewModel};
\node[box, right=of vm] (svc) {GeminiAiService\\(Singleton)};
\node[cloud, right=of svc] (or) {OpenRouter\\API};
\node[box, below=of or] (model) {LLM\\(Llama / Gemma)};
\node[box, below=of svc] (fb) {Firestore\\(resultado)};

\draw[arr] (vm) -- node[above,font=\tiny]{analyzeReport()} (svc);
\draw[arr] (svc) -- node[above,font=\tiny]{HTTP POST} (or);
\draw[arr] (or) -- node[right,font=\tiny]{routing} (model);
\draw[arr] (model) -- node[right,font=\tiny]{JSON} (or);
\draw[arr] (or) -- node[above,font=\tiny]{respuesta} (svc);
\draw[arr] (svc) -- node[right,font=\tiny]{AiReportAnalysis} (fb);
\draw[arr] (fb) -- node[below,font=\tiny]{UI update} (vm);

\end{tikzpicture}
\end{center}

\subsection{Implementación técnica}

La comunicación con OpenRouter se realiza mediante \texttt{java.net.HttpURLConnection}, sin dependencias externas adicionales. La clase \texttt{GeminiAiService} gestiona:

\begin{itemize}[leftmargin=1.5em]
    \item \textbf{Fallback multi-modelo}: si un modelo retorna error 429 (rate limit) o 404, el sistema espera el tiempo indicado por el servidor y pasa automáticamente al siguiente modelo de la lista.
    \item \textbf{Análisis heurístico local}: si todos los modelos fallan, se aplica un análisis basado en reglas locales (longitud del texto, coherencia léxica, categoría) para no bloquear al administrador.
    \item \textbf{Parser JSON robusto}: extracción del objeto JSON mediante balance de llaves, resistente a respuestas malformadas del modelo.
\end{itemize}

\begin{lstlisting}[language=Kotlin, caption={Prompt de moderación enviado al LLM}]
// Sistema:
"Eres un moderador estricto de reportes ciudadanos. Responde ÚNICAMENTE con JSON."

// Usuario:
"""
Título: "${report.title}"
Categoría: ${report.type}
Descripción: "${report.description}"

CAUSAS DE RECHAZO (basta con UNA):
- Texto sin sentido, aleatorio o incomprensible
- Descripción vacía o demasiado corta
- No relacionado con vías, infraestructura o seguridad vial
- Contenido ofensivo o spam

Responde con: {"recommendation":"VERIFICAR"|"RECHAZAR",
               "confidence":"Alta"|"Media"|"Baja",
               "keyPoints":[...], "reasoning":"..."}
"""
\end{lstlisting}

\subsection{Claves API}

Las claves de API se almacenan en \texttt{local.properties} (excluido del control de versiones) y se inyectan en tiempo de compilación mediante \texttt{BuildConfig}, garantizando que no queden expuestas en el repositorio.

% ════════════════════════════════════════════
\section{Notificaciones Push: Firebase Cloud Messaging}
% ════════════════════════════════════════════

\subsection{Descripción}

Las notificaciones push se implementaron con \textbf{Firebase Cloud Messaging (FCM)}, permitiendo informar a los ciudadanos sobre cambios en el estado de sus reportes en tiempo real, sin necesidad de que la aplicación esté activa.

\subsection{Flujo de notificaciones}

\begin{enumerate}[leftmargin=1.5em]
    \item El administrador \textbf{verifica} o \textbf{rechaza} un reporte desde el panel de moderación.
    \item El \texttt{AdminReportModerationViewModel} actualiza el estado en Firestore y llama a \texttt{sendNotificationToUser()}.
    \item La función Cloud guarda una notificación en el documento del usuario en Firestore con el título, mensaje, tipo de acción y razón de rechazo (si aplica).
    \item El dispositivo del ciudadano recibe la notificación a través del servicio \texttt{VialertFirebaseMessagingService}.
\end{enumerate}

\subsection{Implementación del servicio}

\begin{lstlisting}[language=Kotlin, caption={VialertFirebaseMessagingService.kt}]
class VialertFirebaseMessagingService : FirebaseMessagingService() {

    override fun onMessageReceived(message: RemoteMessage) {
        val title = message.notification?.title ?: message.data["title"] ?: "Vialert"
        val body  = message.notification?.body  ?: message.data["body"]  ?: "Nueva notificación"
        showNotification(title, body)
    }

    private fun showNotification(title: String, body: String) {
        val channel = NotificationChannel(CHANNEL_ID, CHANNEL_NAME,
                                          NotificationManager.IMPORTANCE_HIGH)
        notificationManager.createNotificationChannel(channel)

        val notification = NotificationCompat.Builder(this, CHANNEL_ID)
            .setSmallIcon(R.mipmap.vialert)
            .setContentTitle(title)
            .setContentText(body)
            .setStyle(NotificationCompat.BigTextStyle().bigText(body))
            .setPriority(NotificationCompat.PRIORITY_HIGH)
            .setAutoCancel(true)
            .build()

        notificationManager.notify(System.currentTimeMillis().toInt(), notification)
    }
}
\end{lstlisting}

El servicio está registrado en \texttt{AndroidManifest.xml} con el filtro de intención \texttt{com.google.firebase.MESSAGING\_EVENT}.

\subsection{Tipos de notificaciones}

\begin{center}
\begin{tabularx}{0.9\textwidth}{lX}
\toprule
\textbf{Evento} & \textbf{Mensaje enviado al ciudadano} \\
\midrule
Reporte verificado & ``Tu reporte \textit{[título]} ha sido verificado. Gracias por contribuir.'' \\
Reporte rechazado & ``Tu reporte \textit{[título]} fue rechazado. Razón: \textit{[motivo]}.'' \\
\bottomrule
\end{tabularx}
\end{center}

% ════════════════════════════════════════════
\section{Distribución: Firebase App Distribution}
% ════════════════════════════════════════════

El APK de prueba se publica y distribuye utilizando \textbf{Firebase App Distribution}, accesible desde la consola de Firebase en \url{https://console.firebase.google.com}. Esta herramienta permite distribuir versiones de prueba a evaluadores sin necesidad de publicar en Google Play Store.

\textbf{Pasos realizados para la distribución}:
\begin{enumerate}[leftmargin=1.5em]
    \item Generación del APK de release desde Android Studio: \textit{Build → Generate Signed APK}.
    \item Acceso a la sección \textit{App Distribution} en la consola de Firebase del proyecto \textbf{Vialert}.
    \item Carga del APK y configuración de la lista de evaluadores (correos institucionales).
    \item Envío de invitaciones de instalación a los evaluadores.
\end{enumerate}

% ════════════════════════════════════════════
\section{Resumen de Requerimientos Implementados}
% ════════════════════════════════════════════

\begin{center}
\begin{tabularx}{\textwidth}{lXl}
\toprule
\textbf{\#} & \textbf{Requerimiento} & \textbf{Estado} \\
\midrule
1 & Persistencia con Firebase Firestore + caché offline & \textcolor{green!60!black}{✓ Cumplido} \\
2 & Autenticación con Firebase Auth & \textcolor{green!60!black}{✓ Cumplido} \\
3 & Recuperación de contraseña con Firebase & \textcolor{green!60!black}{✓ Cumplido} \\
4 & Funcionalidad de mapas (Mapbox Maps SDK v11) & \textcolor{green!60!black}{✓ Cumplido} \\
5 & Carga de imágenes en la nube (Cloudinary) & \textcolor{green!60!black}{✓ Cumplido} \\
6 & Integración de IA (OpenRouter + LLMs gratuitos) & \textcolor{green!60!black}{✓ Cumplido} \\
7 & Notificaciones push (Firebase Cloud Messaging) & \textcolor{green!60!black}{✓ Cumplido} \\
8 & Distribución APK (Firebase App Distribution) & \textcolor{green!60!black}{✓ Cumplido} \\
\bottomrule
\end{tabularx}
\end{center}

% ════════════════════════════════════════════
\section{Evidencias Visuales}
% ════════════════════════════════════════════

Las evidencias visuales del funcionamiento de las principales características de la aplicación se presentan en el video adjunto a esta entrega. El video incluye demostración de:

\begin{itemize}[leftmargin=1.5em]
    \item Registro e inicio de sesión con Firebase Auth
    \item Recuperación de contraseña
    \item Creación de reporte con selección de ubicación en Mapbox y carga de foto a Cloudinary
    \item Feed de reportes con filtro por distancia GPS
    \item Mapa interactivo de reportes (pantalla Explorar)
    \item Análisis automático de reportes con IA (moderación)
    \item Asistente virtual conversacional
    \item Notificación push al verificar / rechazar un reporte
    \item Sistema de niveles de reputación ciudadana
    \item Panel de administración con estadísticas
\end{itemize}

\vfill
\begin{center}
    \textcolor{gray}{\small Documento generado para la entrega final de Proyecto de Desarrollo Móvil --- Universidad del Quindío --- Mayo 2026}
\end{center}

\end{document}
